%!TEX TS-program = xelatex
% Canonical CV. Single-column layout; role variants retain their own templates.
\documentclass[11pt,a4paper]{article}
\usepackage[margin=19mm,top=17mm,bottom=18mm]{geometry}
\usepackage{fontspec}
\setmainfont[BoldFont=texgyreheros-bold.otf,ItalicFont=texgyreheros-italic.otf]{texgyreheros-regular.otf}
\usepackage{xcolor}
\definecolor{accent}{HTML}{1E6291}
\definecolor{muted}{HTML}{505860}
\usepackage{hyperref}
\hypersetup{colorlinks=true,urlcolor=accent,linkcolor=accent,
  pdfauthor={Christos Hadjinikolis},
  pdftitle={Christos Hadjinikolis | Engineering Manager}}
\usepackage{enumitem}
\usepackage{titlesec}
\usepackage{fancyhdr}
\setlength{\parindent}{0pt}
\setlength{\parskip}{5pt}
\setlist[itemize]{leftmargin=1.2em,itemsep=5pt,topsep=4pt,parsep=0pt}
\titleformat{\section}{\large\bfseries\color{accent}}{}{0pt}{}
\titlespacing*{\section}{0pt}{12pt}{5pt}
\pagestyle{fancy}
\fancyhf{}
\fancyfoot[L]{\footnotesize\color{muted}Christos Hadjinikolis | Engineering Manager}
\fancyfoot[R]{\footnotesize\color{muted}\thepage/2}
\renewcommand{\headrulewidth}{0pt}
\setlength{\footskip}{22pt}
\emergencystretch=2em
\raggedright
\newcommand{\role}[3]{\textbf{#1} \hfill {\small #2}\\{\color{muted}#3}\par}
\begin{document}
{\LARGE\bfseries Christos Hadjinikolis}\par
{\large Engineering Manager | Production ML \& Data Platforms}\par
\href{mailto:christos.hadjinikolis@gmail.com}{christos.hadjinikolis@gmail.com}\\
\href{https://christos-hadjinikolis.github.io/}{Website} \enspace|\enspace
\href{https://uk.linkedin.com/in/christoshadjinikolis}{LinkedIn} \enspace|\enspace
\href{https://github.com/Christos-Hadjinikolis}{GitHub}

\section{Profile}
Engineering Manager leading teams that turn research and complex data into reliable products. Manage six direct reports and lead a ten-person cross-functional team at Vortexa, with responsibility for people development, technical direction and delivery. Engineering Manager since February 2022, following industry roles in ML engineering and consulting since 2016 and earlier doctoral research and teaching. Combine hands-on systems judgement with coaching, stakeholder alignment and production ownership.

\section{Experience}
\role{Vortexa Ltd, London}{December 2020--present}{Engineering Manager | ML Systems Lead (February 2022--present)\\Previously Senior ML Engineer (December 2020--February 2022)}
Lead engineering strategy and delivery for live ML and data services supporting maritime intelligence. The platform processes approximately \textbf{6 million vessel-position records/hour} for \textbf{13,500 monitored vessels}.
\begin{itemize}
  \item \textbf{People leadership.} Manage six direct reports across Data Science, Data Engineering, Product and domain analysis; lead a ten-person cross-functional delivery team. Own performance reviews, career development and delivery priorities. Retained the full team and supported every member's promotion or progression (July 2026 team snapshot).
  \item \textbf{Hiring and organisational growth.} Hiring manager for six roles over time. Shaped hiring practices, system-design interviews, onboarding and mentoring as Data Production grew from four to more than thirty people.
  \item \textbf{Delivery improvement.} Introduced local end-to-end tests and reusable data-access patterns, cutting a three-hour pipeline development feedback loop to under five minutes. Reduced reliance on individual reviewers by making changes easier for the team to test and understand.
  \item \textbf{Platform direction.} Led the Kafka Streams-to-Flink migration for live pipelines, enabling compute scaling independent of Kafka partitioning and improving operational visibility and maintainability. Established shared on-call, runbooks and rollback/fallback practices.
  \item \textbf{Research to production.} Led delivery of destination and arrival-time sequence/transformer models in PyTorch. Established model versioning, deployment, evaluation and monitoring practices so the team could investigate failures and improve production predictions with domain experts.
  \item \textbf{Product and engineering alignment.} Led workshops to resolve competing definitions of model quality and turn them into measurable objectives and prioritised work. Started a weekly architecture forum to share decisions and reduce knowledge silos across teams.
\end{itemize}

\newpage
\section{Earlier Experience}
\role{Data Reply, London}{April 2016--December 2020}{Senior Consultant | Data Scientist / ML Engineer}
Joined the London spin-off as its first consultant and contributed to a practice that grew to more than thirty consultants. Progressed to Senior Consultant, combining hands-on delivery with client scoping, technical leadership, mentoring and consultant placement.
\begin{itemize}
  \item \textbf{Vodafone:} Led technical workstreams on a GCP/Kubeflow data-science platform, covering migration planning, CI/CD and ML workflows; developed a feature-engineering framework for mobile-network analytics.
  \item \textbf{CNHi:} Lead Data Scientist/Scrum Master for predictive maintenance using live vehicle telemetry. Guided a PySpark DS/DE team and translated stakeholder requirements into a delivery backlog.
  \item \textbf{UBS:} Delivered graph/process-mining analytics, observability and Kafka/Elasticsearch pipelines as an embedded consultant; developed production engineering practice through XP and pair programming.
\end{itemize}
\role{KCL, UCL and David Game College}{2010--2016}{Teaching and Technical Training}
Taught algorithms, data structures, databases and Java/Python during and after doctoral research; developed technical communication and mentoring practice.

\section{Technical Experience}
\textbf{Production ML:} PyTorch, MLflow, model serving, evaluation and monitoring.\\
\textbf{Data and platforms:} Python, Java/Kotlin, SQL, Kafka, Flink, AWS/GCP.\\
\textbf{Engineering practice:} CI/CD, data contracts, integration testing, observability and architecture reviews.

\section{Selected Engineering Work}
\textbf{\href{https://pypi.org/project/dynamicio/}{dynamicio}} --- Created a Python library for explicit data-access boundaries and schema validation, used across 15+ repositories to support local testing and reusable ML/data workflows.

\textbf{\href{https://pypi.org/project/skeleton-replay/}{skeleton-replay} and \href{https://plugins.jetbrains.com/plugin/32807-skeleton-replay}{JetBrains plugin}} --- Published developer tools that turn Python runs into traces, architecture views and replayable reports for debugging, review and onboarding.

\textbf{\href{https://christos-hadjinikolis.github.io/2026/08/01/skeleton-replay-runtime-architecture-evidence.html}{Promet} (personal project)} --- Local AI assistant runtime exploring tool execution, approval boundaries, durable state, voice interaction and trace/replay workflows.

\section{Research, Education and Public Work}
\textbf{UCL Associate Researcher} (since October 2024): research and teaching on practical AI standardisation, auditability and responsible adoption.\\
\textbf{AI Standards Committee Expert} (since 2021): ISO/CEN-CENELEC standards work, including JTC 21 WG3.

\textbf{Ph.D. Computer Science}, King's College London, 2010--2014. Persuasion dialogues and opponent modelling; IJCAI 2013 Best Poster Award. \href{https://scholar.google.com/citations?hl=en&user=Rr3KTx8AAAAJ}{Publications}.\\
\textbf{Diploma (BEng), Computer Engineering}, University of Thessaly, 2004--2010. Five-year polytechnic degree; AI specialisation.

\textbf{Selected talks:} Agile in Action (2023), engineering and delivery at Vortexa; Open Data Science Conference (2022), testing and data-access boundaries with dynamicio.
\end{document}
